%!TEX program = xelatex
\documentclass[11pt,a4paper]{article}
\usepackage[margin=2.2cm]{geometry}
\usepackage{amsmath,amssymb,amsthm,mathtools}
\usepackage{bm}
\usepackage{enumitem}
\usepackage{booktabs}
\usepackage{xcolor}
\usepackage{hyperref}
\usepackage{fancyhdr}
\usepackage{xeCJK}

\hypersetup{
    colorlinks=true,
    linkcolor=blue!60!black,
    citecolor=blue!60!black,
    urlcolor=blue!60!black
}

\pagestyle{fancy}
\fancyhf{}
\rhead{\small Qiuzhen QE Probability \& Statistics --- Fall 2024}
\lhead{\small Official Qualifying Exam}
\rfoot{\small Page \thepage}

\DeclareMathOperator*{\argmax}{arg\,max}
\DeclareMathOperator*{\argmin}{arg\,min}
\DeclareMathOperator{\Var}{Var}
\DeclareMathOperator{\E}{\mathbb{E}}
\DeclareMathOperator{\Pbb}{\mathbb{P}}
\DeclareMathOperator{\Cov}{Cov}

\setlist[itemize]{topsep=3pt,itemsep=2pt,parsep=0pt}
\setlist[enumerate]{topsep=3pt,itemsep=2pt,parsep=0pt}

\title{\textbf{QUALIFYING EXAM: PROBABILITY \& STATISTICS}\\[0.5em]\Large FALL 2024}
\date{}
\author{}

\begin{document}

\maketitle
\vspace{-3em}

\section*{Instructions}
\begin{itemize}
    \item There are 11 problems in this exam. You need to choose 8 of them to solve. If you select more than 8, only the first 8 that you have worked on will be graded. Note that 4 of the problems are worth 15 points each and the rest 10 points each.
    \item You must follow all the rules of exam taking. Misconducts will be subject to proper disciplinary actions by the Center.
    \item You must provide all necessary details for full credits. A final answer with no or little explanation/derivation, even if correct, receives a minimal credit.
    \item $\mathbb{R}$ denotes the set of real numbers and $\mathbb{N} = \{1, 2, 3, \dots\}$ denotes the set of positive integers. $\xrightarrow{(d)}$ and $\stackrel{(d)}{=}$ mean ``converges in distribution'' and ``equal in distribution'', respectively.
\end{itemize}

\noindent\rule{\linewidth}{0.4pt}

\begin{enumerate}[label=\textbf{\arabic*.}]
    \item \textbf{(10 points)} Let $U_1, U_2, \dots$ be independent identically distributed (i.i.d.) random variables uniformly distributed on $[0, 1]$, and define $S_n = \sum_{k=1}^n U_k$.
    \begin{enumerate}[label=(\alph*)]
        \item Calculate $\E[(S_n)^4]$.
        \item Determine the distribution of $S_3$.
    \end{enumerate}

    \item \textbf{(10 points)} Let $(Y_n)_{n \ge 1}$ be a sequence of real-valued random variables such that
    \[
    \sqrt{n}(Y_n - a) \xrightarrow[n \to \infty]{(d)} \mathcal{N}(0, \sigma^2),
    \]
    where $\mathcal{N}(0, \sigma^2)$ stands for the normal distribution with $\sigma \neq 0$.
    \begin{enumerate}[label=(\alph*)]
        \item Let $g : \mathbb{R} \to \mathbb{R}$ be a function differentiable at $a$ with $g'(a) \neq 0$. Prove that
        \[
        \sqrt{n}(g(Y_n) - g(a)) \xrightarrow[n \to \infty]{(d)} \mathcal{N}\left(0, [g'(a)]^2 \sigma^2\right).
        \]
        \item Fix $p \in (0, 1)$. For $n \in \mathbb{N}$, let $Z_n \stackrel{(d)}{=} \mathrm{Bin}(n, p)$ be a binomial random variable. Prove that
        \[
        \sqrt{n} \left( \ln\left(\frac{Z_n}{n}\right) - \ln p \right) \xrightarrow[n \to \infty]{(d)} \mathcal{N}\left(0, \frac{1 - p}{p}\right).
        \]
    \end{enumerate}

    \item \textbf{(10 points)} Let $X$ be a real-valued random variable. Prove that the following properties are equivalent in the sense that the parameters $K_i > 0$ appearing in these properties differ from each other by at most an absolute constant factor:
    \begin{enumerate}[label=(\alph*)]
        \item The tails of $X$ satisfy $\Pbb(|X| \ge t) \le 2 \exp(-t / K_1)$ for all $t \ge 0$.
        \item The moments of $X$ satisfy $\E(|X|^p) \le (K_2 p)^p$ for all $p \in \mathbb{N}$.
        \item The moment generating function of $|X|$ is bounded at some point, i.e., $\E \exp(|X| / K_3) \le 2$ for some $K_3 > 0$.
    \end{enumerate}
    (\textit{Hint:} You can use Stirling's approximation: $e(n/e)^n \le n! \le e n (n/e)^n$ for all $n \in \mathbb{N}$.)

    \item \textbf{(10 points)} Let $\mathbb{T}_d$ be an infinite $d$-regular tree, where every node has degree $d$. On the other hand, let $\mathcal{T}_b$ be an infinite $b$-ary tree with root $o$. Consider simple random walk $X_n$ on $\mathbb{T}_d$ and simple random walk $Y_n$ on $\mathcal{T}_b$.
    \begin{enumerate}[label=(\alph*)]
        \item For each $d \in \mathbb{N}$ with $d \ge 2$, determine whether $X_n$ on $\mathbb{T}_d$ is recurrent or transient. Prove your claim.
        \item For each $b \in \mathbb{N}$, determine whether $Y_n$ on $\mathcal{T}_b$ is recurrent or transient. Prove your claim.
    \end{enumerate}

    \item \textbf{(15 points)} Let $X_1, X_2, \dots$ be i.i.d.\ random variables with exponential distribution: $\Pbb(X_k > x) = e^{-x}$ for $x \ge 0$. Define $M_n \coloneqq \sum_{k=1}^n \frac{X_k}{k}$.
    \begin{enumerate}[label=(\alph*)]
        \item Prove that $(M_n - \ln n)_{n \in \mathbb{N}}$ converges to a limit $Y$ almost surely.
        \item Prove that, for every $p \in (0, 1)$, $\left( \frac{\exp(p M_n)}{n^p} \right)_{n \in \mathbb{N}}$ converges to a limit $Z$ in $L^1$.
    \end{enumerate}

    \item \textbf{(15 points)} Let $B_t$ be a 1D standard Brownian motion started from 0.
    \begin{enumerate}[label=(\alph*)]
        \item Consider $X_t = x + B_t$ started at $x > 0$. For any $t > 0$ and $b > a > 0$, compute
        \[
        \Pbb \left( X_t \in [a, b] \;\Big|\; \min_{0 \le s \le t} X_s > 0 \right).
        \]
        \item A Brownian bridge $W_t$ on $[0, 1]$ is $B_t$ conditioned on $B_1 = 0$. A 1D Gaussian free field $h_t$ on $[0, 1]$ with zero boundary has covariance given by the Green's function $G(t, s)$ of $-\Delta$. Prove that $(W_t : t \in [0, 1])$ has the same distribution as $(h_t : t \in [0, 1])$.
    \end{enumerate}

    \item \textbf{(10 points)} Let $X_1, \dots, X_n$ be an i.i.d.\ sample from $\mathcal{N}(\mu, 1)$ with $\mu$ unknown. One only records $\mathbf{Y} = (Y_1, \dots, Y_n)$ where $Y_i = I(X_i < 0)$.
    \begin{enumerate}[label=(\alph*)]
        \item Derive the MLE of $\mu$ based on the observed data $\mathbf{Y}$.
        \item Construct a size $\alpha$ uniformly most powerful (UMP) test for testing $H_0 : \mu \le \mu_0$ versus $H_1 : \mu > \mu_0$ based on $\mathbf{Y}$.
        \item Describe how to construct a $(1 - \alpha)$ confidence interval for $\mu$ based on $\mathbf{Y}$.
    \end{enumerate}

    \item \textbf{(10 points)} Consider the linear model $Y_i = \mathbf{z}_i^T \bm{\beta} + \epsilon_i$, $i = 1, \dots, n$, under Gauss-Markov assumptions ($\E[\epsilon_i] = 0$, $\Var[\epsilon_i] = \sigma^2$, $\Cov(\epsilon_i, \epsilon_j) = 0$ for $i \neq j$).
    \begin{enumerate}[label=(\alph*)]
        \item Let $\hat{\bm{\beta}} = (\mathbf{Z}^T \mathbf{Z})^{-1} \mathbf{Z}^T \mathbf{Y}$ be the least squares estimate. Let $\theta = \mathbf{b}^T \bm{\beta}$. Write down the mean and variance of $\hat{\theta} = \mathbf{b}^T \hat{\bm{\beta}}$, and prove that $\hat{\theta}$ is the best linear unbiased estimator (BLUE) of $\theta$.
        \item Further assume $\epsilon_i \stackrel{\text{i.i.d.}}{\sim} \mathcal{N}(0, \sigma^2)$ with $\sigma^2$ known. Derive the Fisher information matrix $I(\bm{\beta})$.
    \end{enumerate}

    \item \textbf{(10 points)} Suppose $X_1, \dots, X_n$ are i.i.d.\ Uniform$[0, \theta]$ with $\theta > 0$. Fix $t \in (0, \theta)$. Consider $F_n(t) = \frac{1}{n} \sum_{i=1}^n I_{\{X_i \le t\}}$ and $T_n(t) = t / (2\overline{X})$.
    \begin{enumerate}[label=(\alph*)]
        \item Find the asymptotic distributions of the two estimators.
        \item For what value of $t$ will the first estimator have a smaller asymptotic variance than the second estimator?
        \item Let $\theta = 1$. Find the asymptotic distribution of $n F_n(n^{-1/2}) - \sqrt{n}$.
    \end{enumerate}

    \item \textbf{(15 points)} Let $X_1, \dots, X_n$ ($n \ge 2$) be i.i.d.\ $\mathcal{N}(\mu, \sigma^2)$ with $\mu \ge 0, \sigma > 0$ unknown.
    \begin{enumerate}[label=(\alph*)]
        \item Show $\overline{X}$ and $S^2$ are independent.
        \item Find UMVUE of $\mu / \sigma$ if it exists.
        \item Is $\overline{X}$ admissible for estimating $\mu$ under squared error loss? Prove your assertion.
    \end{enumerate}

    \item \textbf{(15 points)} Let $X_1, \dots, X_n$ be an i.i.d.\ sample from Uniform$[\theta, \theta + |\theta|]$ where $\theta \neq 0$.
    \begin{enumerate}[label=(\alph*)]
        \item Derive the method of moments estimator of $\theta$.
        \item Derive the MLE of $\theta$, $\hat{\theta}$.
        \item Is $\hat{\theta}$ a consistent estimator of $\theta$? Explain.
    \end{enumerate}
\end{enumerate}

\end{document}
